%% main.tex -- PLanQC 2027 talk proposal (extended abstract)
%% "Human-Made Scaffolds for Machine-Written Proofs: Sixteen
%% Completeness Theorems for Quantum Circuits in Agda"
%%
%% Derived from the OOPSLA draft in ../oopsla.  PLanQC asks for 2-5
%% pages in single-column 11pt acmsmall, not counting references.
%%
%% Build: make   (latexmk -pdf main.tex)
\documentclass[acmsmall,nonacm,11pt,screen]{acmart}

\settopmatter{printacmref=false, printccs=false, printfolios=true}
\setcopyright{none}
\renewcommand\footnotetextcopyrightpermission[1]{}

%% The authors' names share one line in the title block, each marked
%% by a superscript; the affiliations and e-mail addresses are in the
%% first-page footnote (\authorsaddresses below), under the same marks.
%% The marks stay out of the PDF metadata.
\newcommand{\affmark}[1]{\texorpdfstring{\textsuperscript{#1}}{}}
\makeatletter
%% No final period, which would read as part of the last address.
\def\@setauthorsaddresses{\@authorsaddresses\unskip}
%% Names as written, not in the capitals acmsmall sets them in.
\let\planqc@mkauthors\@mkauthors
\def\@mkauthors{\begingroup\let\MakeUppercase\@firstofone
  \planqc@mkauthors\endgroup}
\makeatother

%% Narrower top and bottom margins than acmsmall's 58pt and 44pt.
\geometry{top=20pt, bottom=24pt}

\usepackage{booktabs}
\usepackage{tabularx}
\newcolumntype{Y}{>{\raggedright\arraybackslash}X}
\usepackage{agda-style}
\usepackage{subcaption}
\usepackage{cleveref}
%% Footnotes below bottom floats, so that the addresses stay at the
%% foot of the first page, under Figure 1.
\usepackage{fnpos}
\makeFNbottom
\makeFNbelow

%% TikZiT figures. A subset of tikzit.sty (the layer stack, the blank
%% node style and \tikzfig), which would load enumitem, which clashes
%% with acmart. proofs.tikzstyles has the styles of figures/*.tikz.
\usepackage{tikz}
\usetikzlibrary{backgrounds,arrows,shapes,shapes.geometric,shapes.misc}
\pgfdeclarelayer{edgelayer}
\pgfdeclarelayer{nodelayer}
\pgfsetlayers{background,edgelayer,nodelayer,main}
\tikzset{tikzit fill/.initial=0,tikzit draw/.initial=0,tikzit shape/.initial=0,tikzit category/.initial=0}
\tikzstyle{none}=[inner sep=0mm]
\tikzstyle{tikzfig}=[baseline=-0.25em,scale=0.7]
\newcommand{\tikzfig}[1]{{\tikzstyle{every picture}=[tikzfig]\input{#1.tikz}}}
\input{proofs.tikzstyles}

\citestyle{acmnumeric}
%% Table 1 keeps author-year citations: ``Greylyn [2014]''.
\newcommand{\tabcite}[1]{\citeauthor{#1}~[\citeyear{#1}]}

%% Inline Agda identifiers, in a sans-serif face.
\newcommand{\aid}[1]{\mbox{\textsf{\small #1}}}
\newcommand{\Z}{\mathbb{Z}}
\newcommand{\D}{\mathbb{D}}

\begin{document}

\title[QCert: Verifying Sixteen Quantum Circuit Completeness Results in One Agda Framework]{QCert: Verifying Sixteen Quantum Circuit Completeness Results in One Agda Framework}
% \subtitle{Extended Abstract}


%Human-Made Scaffolds for Machine-Written Proofs of Completeness
%Sixteen Completeness Theorems for Quantum Circuits in Agda: Human-Guided, Machine-Written Proofs

%% Affiliations and e-mail addresses as in ../vqupit/main.tex; Yuan
%% Feng's address is to be confirmed (see README.md).
\author[Xiaoning Bian]{Xiaoning Bian\affmark{1}}
\author[Yuan Feng]{Yuan Feng\affmark{1}}
\author[Sarah Meng Li]{Sarah Meng Li\affmark{2}}
\author[John van de Wetering]{John van de Wetering\affmark{2}}
\authorsaddresses{%
  \textsuperscript{1}Tsinghua University, Beijing, China,
  \nolinkurl{bian@dal.ca}, \nolinkurl{yuan_feng@tsinghua.edu.cn}\newline
  \textsuperscript{2}University of Amsterdam and QuSoft, Amsterdam,
  The Netherlands, \nolinkurl{sarah.li@uva.nl},
  \nolinkurl{john@vdwetering.name}}

% \begin{abstract}
% % An equational theory for a family of quantum circuits is a set of rewrite rules. The theory 

% A set of circuit equations is complete if any two circuits denoting the same linear map can be derived from each other using these equations, and minimal if no equation can be derived from the others. Establishing these properties takes long case analyses spanning hundreds of pages, making the details hard to check by hand. To this end, we extend the previous formalization framework to QCert, and use it to formalize the main circuit completeness theorems to date. This includes sixteen results, fifteen on completeness and one on minimality, over finite gate sets such as Clifford, CNOT-dihedral, Clifford+T, Clifford+CS, real Clifford+CH, and Toffoli--Hadamard, on qubits and on qudits of every odd prime dimension. QCert divides the work between humans and LLM coding agents. We first built its framework by specifying the syntax, semantics, and group presentations. Then, each paper's main theorem can be expressed uniformly in this framework. Next, an agent took over and formalized the proofs in Agda, following the papers' arguments. Agda checked every proof, so a human need only check the relation table, the gate semantics, and the theorem statement for each paper. Overall, QCert confirms the main results of each work and identifies gaps in two of the existing proofs. Its ability to do verification at scale makes QCert well suited to verifying quantum circuit completeness in higher dimensions.
% \end{abstract}

\maketitle
\vspace{-.4 cm}

\section{Introduction}
An equational theory for a family of quantum circuits is a set of equations between circuits used as rewrite rules. By applying these rules to local subcircuits, we can transform a circuit into an equivalent one. Such a theory should be \emph{sound} because every rewrite preserves the underlying linear map. It should also be \emph{complete} so that any equivalence between circuits can be derived. Additionally, we want it to be \emph{minimal}, with no redundant rules. Equational rewriting supports circuit optimization~\cite{hietala2021voqc,nam2018automated,maslov2008simplification,xu2022quartz}, equivalence checking~\cite{yamashita2010fast}, and hardware mapping~\cite{sivarajah2020tket}. Completeness ensures that every circuit equivalence can be derived within the theory, making it a powerful tool for reasoning about quantum programs. For example,
\mbox{Hadamard-$\Pi$}~\cite{fang2026hadamard} provides a complete
equational framework for reasoning about reversible quantum programs. Since 2014, completeness results have been established for a wide range of circuit families such as Clifford~\cite{selinger2015clifford,makary2021real,li2025qutrit,bian2026qupit}, CNOT-dihedral~\cite{amy2018cnotdihedral,blake2026phaseaffine}, two-qubit Clifford+T~\cite{bian2023cliffordt2}, three-qubit Clifford+CS~\cite{bian2023cliffordcs3}, Toffoli--Hadamard~\cite{amy2024toffoli} and real Clifford+CH~\cite{clement2026realcliffordch} circuits, and with continuous parameters, all quantum circuits~\cite{clement2023complete}. Completeness has also been obtained for the group of two-level matrices underlying several of these circuit families~\cite{greylyn2014,bian2021un,li2021on,fang2026hadamard,clement2019thesis}. Minimal theories have been given for six of these circuit fragments~\cite{blake2026simpler} and for all quantum circuits~\cite{clement2024minimal}.


Completeness proofs use two main strategies. \Cref{fig:normalform} outlines the first one. It specifies a unique normal form for each linear map~\cite{selinger2015clifford,bian2026qupit,makary2021real,amy2018cnotdihedral,blake2026phaseaffine,li2025qutrit}, often through giving an exact synthesis algorithm~\cite{greylyn2014,bian2021un,li2021on,fang2026hadamard,clement2019thesis}. Completeness follows by proving that every circuit can be transformed into its normal form using the rewrite rules. \Cref{fig:translation} illustrates the second strategy, which leverages the complete equational theory of another gate set. Completeness carries over if every rule of that theory, translated back, can be derived from the new rules. The new rules must also prove that each gate is equal to its round-trip translation~\cite{bian2023cliffordt2,bian2023cliffordcs3,clement2026realcliffordch,amy2024toffoli,blake2026simpler,clement2023complete}. Sometimes the two methods are combined, as in the scaled dyadic extension of~\citet{li2021on}. 
%It pairs a normal form with reduction to an already complete subgroup. 

\begin{figure}[H]
\centering
\begin{subfigure}[b]{0.42\linewidth}
  \centering
  \scalebox{.8}{\tikzfig{figures/normalform}}
  \caption{$C_1$ and $C_2$ can be rewritten to the same normal form.}
  \label{fig:normalform}
\end{subfigure}\hfill
\begin{subfigure}[b]{0.52\linewidth}
  \centering
  \scalebox{.8}{\tikzfig{figures/translation}}
  \caption{The complete theory of $\mathcal{G}'$ proves $C_1' = C_2'$. Translating the proof back proves $C_1 = C_2$ over $\mathcal{G}$.}
  \label{fig:translation}
\end{subfigure}
\caption{Two strategies for proving the rewrite rules of a gate set $\mathcal{G}$ are complete. $C_1$ and $C_2$ are any two circuits over $\mathcal{G}$ with the same linear map. In (b), each $C_i$ is translated into a circuit $C_i'$ over another gate set $\mathcal{G}'$ that already has a complete equational theory. The translation preserves the linear map.}
\label{fig:strategies}
\end{figure}

Starting from a complete rule set, we can remove a redundant rule by deriving it from those other rules, thereby preserving
completeness. To show that the resulting set is minimal, we need to give, for each remaining rule, an interpretation of circuits that every other rule preserves except that rule. Any derivation using the other rules preserves this interpretation, so the rule cannot be derived from them~\cite{clement2024minimal,blake2026simpler}.



Completeness and minimality proofs can require long case distinctions
over gates, relations and normal forms, with supporting derivations
sometimes running to hundreds of
pages~\cite{fang2026hadamard,li2025qutrit,clement2026realcliffordch}.
Long reductions of relation sets can also hide circular reasoning,
where a derivation used to eliminate a relation depends indirectly
on that same relation. Such dependencies are hard to detect by hand.
An unnoticed gap can also undermine later proofs that rely on the
result: Section~\ref{sec:findings} reports a gap in a published proof
that a later theorem had imported. A machine check verifies the
individual cases and guards against circular arguments. To check
results that build on one another, it must be done in one framework,
which defines circuits, their semantics and completeness once, so that
a theorem proved in one development can be used in another.


Language models can now write formal proofs of this size; a proof of
Fermat's Last Theorem in Lean, of 13 million lines, was written in 11
days~\cite{anthropic2026flt}. What still needs a person is the
statement: a formal proof proves what its statement says, and someone
must check that the statement, and the definitions under it, mean
what was intended~\cite{hales2026reliability}. Generated proofs are
not written to be read, so that check must stay small. Our approach
is a \emph{scaffold}: code that people write and check once, which
fixes the form of every theorem built on it, so the per-paper check is
short and the rest left to Agda. Its first layer is an Agda library of
circuits and their presentations~\cite{bian2027popl}.

This work grew out of our formalization of the completeness theorem
for qudit Clifford circuits~\cite{bian2026qupit}. There we set up a framework by hand, and brought in LLM coding agents
only at the last stage, to finish the proof. That experience convinced us that all
completeness results for quantum circuits could, and should, be
formalized within one framework. This paper does so: we built such a
scaffold for completeness theorems, and on it formalized sixteen papers
(Table~\ref{tab:corpus}). Our contributions are threefold.
\begin{itemize}
\item QCert: an earlier library of words, circuits and
  presentations~\cite{bian2027popl}, extended with matrix semantics over
  the rings of exact synthesis and a framework for the independence of
  rules, giving every theorem the same form (Section~\ref{sec:form});
\item machine-checked proofs of the sixteen theorems: about 1{,}500
  modules and 420{,}000 lines of Agda with no postulates, checked in
  about 40 minutes (Section~\ref{sec:corpus});
\item a gap in the completeness proof of Hadamard-$\Pi$ and its repair,
  a rule omitted from a Toffoli--Hadamard theory, and misprints in the
  other papers (Section~\ref{sec:findings});
\item a method in which people build the framework and state the
  theorems, LLM coding agents write the proofs, and Agda checks them
  (Section~\ref{sec:method}).
\end{itemize}

\section{One Form for Every Theorem}\label{sec:form}

A relation on an alphabet $X$ generates a congruence $\approx$ on the
words over $X$, and a semantics is a monoid homomorphism
$\llbracket - \rrbracket$ from words to a monoid $M$. It is sound if
related words have equal images and complete if words with equal images
are related; then the relation presents the image of
$\llbracket - \rrbracket$, and if every element of $M$ is an image, it
presents $M$. For circuits, $M$ is a group of matrices, and the image
is the group that the gates generate.

The earlier library~\cite{bian2027popl} fixes this vocabulary: a
presentation theorem is a record, $\Gamma$ \aid{IsSubPresentationOf} $G$,
holding an interpretation of words in a group $G$ with proofs that it
is a group homomorphism on words modulo $\approx$ and injective, and
$\Gamma$ \aid{IsPresentationOf} $G$ adds that it is onto. It also fixes the
syntax of circuits: a gate set is a family indexed by arity, a gate
sits on the bottom wires of a circuit and a shift moves a circuit up a
wire, and any relation stated at the width it is drawn on is extended
by structural rules (an equation holds one wire up, gates on disjoint
wires commute, a scalar is the same on every wire), which are sound
once for every gate set. What this work adds is the semantic side: the
rings of exact synthesis, $\Z[\frac1{\sqrt2},i]$, $\Z[\frac12,i]$,
$\Z[\frac1{\sqrt2}]$ and $\Z[\frac12]$, with their unitary and
orthogonal groups as groups of Agda's standard library, and operators
on bit vectors with the Kronecker calculus, so that the matrix of a
circuit is defined once for every gate set.

Every development ends in one line of this form. The real stabilizer
circuits of~\citet{makary2021real}, with the relations R1--R16 of their
Figure~1 lifted to $n$ wires, conclude
\begin{agdacode}
presentation : {n : ℕ} → (n VRel,_===_) IsSubPresentationOf U (2 ^ n)
\end{agdacode}
and the two-level matrices of Hadamard-$\Pi$, with the relations of
its Figure~6,
\begin{agdacode}
presentation : {n : ℕ} → (_===_ {n}) IsPresentationOf (O n)
\end{agdacode}
To trust such a line a reader checks four things: the relation, as
transcribed from the paper; the group, which comes from the layer and
is checked once for all papers; the interpretation of each gate as a
matrix; and the one line that gives the record that interpretation.
For the real stabilizer development these come to under 500 of its
12{,}010 lines. A single root module, the place to start reading,
states all sixteen main theorems in this form.

\section{Sixteen Formalizations}\label{sec:corpus}

Table~\ref{tab:corpus} lists the sources, what we proved, the size of
each development and its typechecking time. Every theorem is proved
without hypotheses. The two-level presentations hold in every
dimension (Greylyn's for $n\le4$, as in the thesis), and the circuit
theories for every number of wires, except where the paper fixes it.

\begin{table}
\small
\setlength{\tabcolsep}{4pt}
\begin{tabularx}{\linewidth}{@{}l Y l r r@{}}
\toprule
Source & Object & Result & Mod. & Time \\
\midrule
\multicolumn{5}{@{}l}{\emph{Groups of two-level matrices}}\\
\tabcite{greylyn2014} & $U_n(\Z[\frac1{\sqrt2},i])$ & presentation, $n\le4$ & 31 & 42 \\
\tabcite{bian2021un} & $U_n(\Z[\frac12,i])$ & presentation & 37 & 62 \\
\tabcite{li2021on} & $O_n(\Z[\frac12])$, $L_n$ & presentations & 68 & 79 \\
\tabcite{fang2026hadamard} & $O_n(\Z[\frac1{\sqrt2}])$ & presentation & 44 & 44 \\
\tabcite{clement2019thesis} & $O_n(\Z[\frac1{\sqrt2}])$ & presentation & 40 & 64 \\
\midrule
\multicolumn{5}{@{}l}{\emph{Circuits}}\\
\tabcite{selinger2015clifford} & Clifford & sub-presentation & 31 & 34 \\
\tabcite{bian2026qupit} & Clifford, odd prime $p$ & presentation of a model & 215 & 546 \\
\tabcite{makary2021real} & real stabilizer & sub-presentation & 30 & 35 \\
\tabcite{amy2018cnotdihedral} & CNOT-dihedral & sub-presentation & 30 & 24 \\
\tabcite{blake2026phaseaffine} & phase-affine, prime $d$ & sub-presentations & 56 & 61 \\
\tabcite{bian2023cliffordt2} & two-qubit Clifford+T & sub-presentation in $U_4$ & 62 & 118 \\
\tabcite{bian2023cliffordcs3} & three-qubit Clifford+CS & sub-presentation in $U_8$ & 123 & 250 \\
\tabcite{clement2026realcliffordch} & real Clifford+CH & sub-presentation & 346 & 405 \\
\tabcite{amy2024toffoli} & three-qubit Toffoli--Hadamard & presentation of $L_8$ & 201 & 102 \\
\tabcite{li2025qutrit} & qutrit Clifford & presentation & 89 & 282 \\
\midrule
\multicolumn{5}{@{}l}{\emph{Minimality}}\\
\tabcite{blake2026simpler} & six fragments & minimal, complete & 73 & 183 \\
\midrule
All & & & 1{,}476 & 2{,}332 \\
\bottomrule
\end{tabularx}
\captionsetup{position=bottom}
\caption{The object presented, the result (a presentation of the
whole group or of the group the gates generate), the number of modules,
and typechecking time from scratch in seconds, the sum of the modules'
own times given the developments they build on (Agda~2.8, Core
i5-14400).}
\label{tab:corpus}
\end{table}

The papers use few proof methods, and in Agda each becomes a route
with a fixed shape. \emph{Exact synthesis}: an algorithm writes every
matrix as a word, and completeness follows by well-founded induction on
a level from a Main Lemma that closes a square for every state and
generator; the induction is a module written once, and the author
supplies the Main Lemma. \emph{Normal forms}: every circuit is pushed
to a normal form, and the normal form is read back from the matrix.
\emph{Induction on the width}: a circuit is decomposed into a part one
wire down and a part on the bottom wire. \emph{Encoding}: circuits are
translated into a theory already known to be complete and back. And
two ports reuse Bian and Selinger's own Agda
proofs~\cite{bian2023cliffordt2,bian2023cliffordcs3}, which relate the
circuit relations to two-level relations by the Reidemeister--Schreier
method, and connect them to our two-level presentations, so that they
end in a statement about matrices. The qudit Clifford theorem
of~\citet{bian2026qupit} is stated against an algebraic model of the
Clifford group, built from the Pauli and symplectic groups with a
phase; the others are stated against matrices.

Much of each development is generated. The derivations that a paper
gives as computations, such as Selinger's 105-page supplement or the
box relations of the qutrit normal form, come from scripts that follow
the paper or search for the steps, and are checked by a rewriter in
Agda; soundness is decided by evaluating matrices.
What a paper gives as an idea, the Main Lemma or the uniqueness of a
normal form, is written. In the real stabilizer development, half the
lines are generated and a fifth are the paper's mathematics.

\section{What We Found}\label{sec:findings}

\paragraph{A gap in Hadamard-$\Pi$.} Hadamard-$\Pi$ proves that the
relations of its Figure~6 present the group
$O_n(\Z[\frac1{\sqrt2}])$ of orthogonal matrices over
$\Z[\frac1{\sqrt2}]$ (its Theorem~4.14), by the exact synthesis route.
One subcase of its Main Lemma relies on Lemmas~A.16 and~A.17, and both
are false: each fails on a unit column of dimension~6, such as
$(-1-\sqrt2, 1, -1-\sqrt2, 1, \sqrt2, -2+\sqrt2)/\sqrt2^4$, because the
proof drops a $+1$ from a bracket that it reduces modulo~2. Following
the case analysis on random states, we found a state in dimension~9 at
which every diagram the proof offers rises above the level it must stay
below, so the proof of the Main Lemma does not cover it. The theorem is
still true. An unpublished thesis of Cl\'ement from
2019~\cite{clement2019thesis} presents the same group by other
relations, with a case that covers exactly this edge; we proved the two
relation sets equivalent in every dimension, formalized his case as a
decision tree that Agda checks by evaluation, and then his whole proof.
Cl\'ement's 2026 theorem on real Clifford+CH
circuits~\cite{clement2026realcliffordch} imports Theorem~4.14, so it
rested on the same gap; our formalization of it discharges the import.

\paragraph{An omitted relation.} Amy, Ross and
Wesley~\cite{amy2024toffoli} present three-qubit Toffoli--Hadamard
circuits in two steps, first $O_8(\Z[\frac12])$ by the Toffoli and
$\mathrm{K} = \mathrm{H}\otimes\mathrm{H}$ gates, then the scaled dyadic orthogonal group $L_8$.
Their Section~5.5 uses a relation $r_{\mathrm{K}}$ that defines $\mathrm{K}_{1,2}$ by a
controlled K, but the figures that define the relations of their
Theorem~5.2 omit it, and it does not follow from them: reading each
Toffoli-K gate $g$ as $RgR$, with $R = I - J/4$ for the all-ones
matrix $J$, satisfies every other relation and not $r_{\mathrm{K}}$. So the
relations as printed are not complete; with $r_{\mathrm{K}}$ restored, we prove
that they present $L_8$.

\paragraph{Misprints.} The other papers are correct in what we
formalized, up to misprints, several of them in statements that proofs
use: eleven of the 137 typed relations of~\citet{makary2021real}, one
relation of the $L_n$ table of~\citet{li2021on}, the scalar of a qutrit
rule in~\citet{blake2026simpler}, two equations of the encoding
of~\citet{clement2026realcliffordch}, five formulas
of~\citet{amy2024toffoli}, and one box relation of~\citet{li2025qutrit}.
Each was found because the statement failed a numerical check or a
derivation failed at the step that used it.

\section{How the Proofs Were Produced}\label{sec:method}

Besides the qudit Clifford development, LLM coding agents wrote the
other fifteen under our direction, within four weeks, most of them in
less than a day each. Three habits made this possible. Every claim is
checked numerically before it is formalized, with exact matrices in a
short script, which caught most misprints before any Agda was written.
Long derivations are generated, by scripts that follow the paper's own
steps, so that a failing step points at a line of the paper. And the scaffold leaves the agent no choice about
the statement: the relation, the group and the record are fixed, and
only the proof remains.

The main cost was Agda's, not the agent's. The conversion checker
compares terms by unfolding them, so two readings of a long product of
matrices that are not syntactically equal can take exponential time;
readings of circuits are therefore abstract, and equations between
matrices are decided as one boolean computed by row operations, a
hundred to four hundred times faster than comparing entry by entry. Everything is constructive, with no excluded middle,
choice or real numbers, so every case split is a decision procedure
and every proof is a program: with soundness, completeness and
decidable matrix equality, each theory's word problem is decided by
running the proof.

% \section{The Talk}

% The talk will show the form of a theorem and what a reader checks, one
% development from its gate set to its theorem, the gap in
% Hadamard-$\Pi$ and how it was found, and what working with coding
% agents on a human-made scaffold changes for proofs of this kind: the
% cost moves from writing proofs to checking statements, and a scaffold
% keeps that check small.

\section{Conclusion and Future Work}\label{sec:conclusion}

All sixteen developments share one framework, and this matters in
three ways. First, every theory, whatever its gate set or proof
method, is stated with the same circuits, presentation records and
matrix groups over the same rings: a common language in which rewrite
rules for different circuit families can be related, as when we proved Cl\'ement's
relations and Hadamard-$\Pi$'s equivalent. Second, it lays the
foundation for completeness in higher-dimensional systems, whose case
analyses outgrow what a reader can check (the qutrit theorem rests on
hundreds of box relations): soundness is decided by evaluating exact
matrices, derivations are generated and machine-checked, and the corpus typechecks in about 40 minutes. Third,
it paves the way for number-theoretic exact synthesis: its rings, their
unitary and orthogonal groups and the route from a Main Lemma to
completeness (\Cref{sec:corpus}) are formalized once, so a new
synthesis algorithm can be verified efficiently, including its Main
Lemma, where the gap in Hadamard-$\Pi$ lay (\Cref{sec:findings}).

\paragraph{Future work.} We are porting the developments to Lean, which
has a much larger community and whose library Mathlib covers far more
of mathematics, so that our theorems can be used alongside it. Going
forward, we will use this framework to solve new problems in circuit
completeness and exact synthesis, such as the Clifford completeness in all dimensions and optimising T-count inmulti-qubit Clifford+T exact synthesis.
%We are also working on a longer version of this work, with the details of each development.

\bibliographystyle{ACM-Reference-Format}
\bibliography{refs}

\end{document}
